\documentclass[11pt]{article}

\usepackage[T1]{fontenc}
\usepackage[utf8]{inputenc}
\usepackage{lmodern}
\usepackage[a4paper,margin=27mm]{geometry}
\usepackage{amsmath,amssymb,amsthm,mathtools}
\usepackage{microtype}
\usepackage{xcolor}
\usepackage{enumitem}
\usepackage{booktabs}
\usepackage{tabularx}
\usepackage{hyperref}

\hypersetup{
  colorlinks=true,
  linkcolor=blue!55!black,
  citecolor=blue!55!black,
  urlcolor=blue!55!black
}

\newtheorem{theorem}{Theorem}[section]
\newtheorem{proposition}[theorem]{Proposition}
\newtheorem{lemma}[theorem]{Lemma}
\newtheorem{corollary}[theorem]{Corollary}
\newtheorem{definition}[theorem]{Definition}
\theoremstyle{remark}
\newtheorem{remark}[theorem]{Remark}

\newcommand{\Hh}{\mathcal H}
\newcommand{\F}{\mathcal F}
\newcommand{\C}{\mathbb C}
\newcommand{\N}{\mathbb N}
\newcommand{\id}{\mathbf 1}
\newcommand{\Tr}{\operatorname{Tr}}
\newcommand{\rank}{\operatorname{rank}}
\newcommand{\gap}{\operatorname{gap}}
\newcommand{\supp}{\operatorname{supp}}
\newcommand{\norm}[1]{\left\lVert #1\right\rVert}
\newcommand{\ip}[2]{\left\langle #1,#2\right\rangle}

\title{The Volume-Independent DET16 Mirror-Gap Theorem:\\
Exact Product Gap, Spin(10) Symmetry, and the Honest Stability Threshold}
\author{Proof memorandum}
\date{31 August 2026}

\begin{document}
\maketitle

\begin{abstract}
On every finite set of disjoint sixteen-mode clusters we prove exactly
that the DET16 Hamiltonian
\[
 H_0=\sum_x(\id-P_{\phi,x})
\]
has the unique product ground state
$\bigotimes_x(|0\rangle_x+e^{i\phi}|F\rangle_x)/\sqrt2$, spectrum
$\{0,1,\ldots,|\Lambda|\}$, and gap one independently of volume.
The proof includes the graded-CAR argument: the determinant monomial has
even degree sixteen, so local projectors on disjoint clusters commute
without an ordering sign.  We also prove local Spin(10) Gauss invariance,
fermion parity, translation symmetry, absence of spontaneous symmetry
breaking, and absence of torus degeneracy.

For a bounded finite-range even perturbation, the
Michalakis--Zwolak stability theorem applies because frustration
freeness, Local-TQO, and the local gap are exact, with
$\Delta_0(\ell)=0$ and $\gamma(r)=1$.  It gives a
volume-independent threshold $t_*>0$ and perturbed gap at least $1/2$.
For finite $\mathbb Z_2/\mathbb Z_4$ dynamical links on an open chain,
the Fröhlich--Pizzo Lie--Schwinger theorem applies after an exact
ambient-cell and endpoint repair.  It gives a unique physical ground
state and gap at least $\Delta/2$, uniformly in chain length, for
$|t|<\Delta\tau_{\rm FP}/m$, where
$\Delta=\min(1,h_E)$ and $\tau_{\rm FP}>0$ is only partially explicit
in the cited proof.  It does \emph{not} give the proposed numerical
threshold $1/(2\cdot16z)$ for ordinary hopping.  The numeric twins pass
all checks and find no finite-volume closing in their declared scans.
The final box now isolates compact-group rotor links, wall mixing, and
the three-spatial-dimensional placement; none is included in the
finite-group chain theorem.
\end{abstract}

\tableofcontents

\section{Concrete model and frozen finite input}

\subsection{One DET16 cluster}

Let $K\cong\C^{16}$ carry the positive-chirality spin representation
$\rho_{16}$ of $\operatorname{Spin}(10)$ and let
$\F(K)=\bigoplus_{n=0}^{16}\bigwedge^nK$ be its fermionic Fock space.
For an ordered orthonormal basis $(e_1,\ldots,e_{16})$, write
\[
 \Omega^\dagger
 =a^\dagger(e_1)\cdots a^\dagger(e_{16}),\qquad
 |F\rangle=\Omega^\dagger|0\rangle ,
\]
and $\Omega=(\Omega^\dagger)^*$.  Put
\[
\begin{aligned}
 A_\phi
 &=P_0+P_{16}+e^{-i\phi}\Omega+e^{i\phi}\Omega^\dagger,\\
 |\Omega_\phi\rangle
 &=\frac{|0\rangle+e^{i\phi}|F\rangle}{\sqrt2},\\
 P_\phi&=\frac12A_\phi
 =|\Omega_\phi\rangle\langle\Omega_\phi|,\qquad
 q_\phi=\id-P_\phi.
\end{aligned}
\tag{1.1}\label{eq:cluster}
\]
The exact finite computation in
\texttt{verification/v1002\_spin10\_mirror\_projector.py} verifies
\[
 \Omega^\dagger\Omega=P_{16},\quad
 \Omega\Omega^\dagger=P_0,\quad
 P_\phi^2=P_\phi=P_\phi^*,\quad
 \rank P_\phi=1,
\tag{1.2}\label{eq:frozen}
\]
and all forty-five $\mathfrak{so}(10)$ commutators vanish.  These are
inputs to this memorandum; the algebraic consequences needed below
are also proved directly.

\subsection{Many disjoint clusters}

Let $\Lambda$ be any nonempty finite set of clusters.  Fix an ordering
only to identify the global Fock space with an ordinary tensor product,
\[
 \F\!\left(\bigoplus_{x\in\Lambda}K_x\right)
 \cong\widehat{\bigotimes}_{x\in\Lambda}\F(K_x),
\tag{1.3}\label{eq:gradedtensor}
\]
where the hat recalls that the intrinsic product is graded.  Operators
have fermion parity $|A|\in\{0,1\}$.  If $A_x$ and $B_y$ are homogeneous
and supported on disjoint clusters, the CAR give
\[
 A_xB_y=(-1)^{|A_x||B_y|}B_yA_x.
\tag{1.4}\label{eq:gradedcommute}
\]
The terms $P_{0,x}$ and $P_{16,x}$ preserve particle number, while
$\Omega_x$ and $\Omega_x^\dagger$ contain sixteen odd generators.
Every summand of $P_{\phi,x}$ is therefore even.  Hence
\[
 [P_{\phi,x},P_{\phi,y}]=[q_{\phi,x},q_{\phi,y}]=0
 \qquad(x\ne y)
\tag{1.5}\label{eq:commute}
\]
with no ordering-dependent sign.  Define
\[
 H_{0,\Lambda}=\sum_{x\in\Lambda}q_{\phi,x}.
\tag{1.6}\label{eq:H0}
\]
Unless stated otherwise, the same phase $\phi$ is used at every site.

\section{Exact spectrum and the volume-independent unit gap}

\begin{lemma}[Order independence of even local products]
\label{lem:order}
For any subset $S\subseteq\Lambda$, the product
$\prod_{x\in S}P_{\phi,x}$ is independent of the order chosen for its
factors.  Under any ordered tensor-product identification
\eqref{eq:gradedtensor}, it is the ordinary tensor product of the local
rank-one projectors on $S$ and the identity on $S^c$.
\end{lemma}

\begin{proof}
Each $P_{\phi,x}$ is even.  Equation \eqref{eq:gradedcommute} therefore
has sign $(-1)^{0\cdot0}=1$ for every transposition of two factors.
Any two orderings differ by transpositions, proving the first claim.
An even operator crosses all preceding graded tensor factors without a
Jordan--Wigner parity string.  Consequently its representative is the
ordinary local matrix on its own factor and identities elsewhere.
Multiplying these representatives proves the second claim.
\end{proof}

\begin{theorem}[T1: exact DET16 product spectrum]
\label{thm:T1}
Let $N=|\Lambda|$ and $d=\dim\F(\C^{16})=2^{16}$.  The commuting
Hamiltonian \eqref{eq:H0} has
\[
 \operatorname{spec}(H_{0,\Lambda})=\{0,1,\ldots,N\}.
\tag{2.1}\label{eq:spectrum}
\]
The multiplicity of $k$ is
\[
 \binom Nk(d-1)^k.
\tag{2.2}\label{eq:multiplicity}
\]
Its unique normalized ground state is
\[
 |\Omega_{\Lambda,\phi}\rangle
 =\widehat{\bigotimes}_{x\in\Lambda}|\Omega_\phi\rangle_x,
\tag{2.3}\label{eq:productground}
\]
and
\[
 \gap(H_{0,\Lambda})=1
\tag{2.4}\label{eq:gapone}
\]
for every finite volume and every boundary condition.
\end{theorem}

\begin{proof}
On one cluster,
$\F(K_x)=\operatorname{Ran}P_{\phi,x}\oplus
\operatorname{Ran}q_{\phi,x}$ with dimensions $1$ and $d-1$.
By Lemma~\ref{lem:order}, the global Hilbert space is the orthogonal
direct sum
\[
 \bigoplus_{S\subseteq\Lambda}
 \left(
 \bigotimes_{x\in S}\operatorname{Ran}q_{\phi,x}
 \otimes
 \bigotimes_{y\notin S}\operatorname{Ran}P_{\phi,y}
 \right).
\tag{2.5}\label{eq:syndrome}
\]
On the summand labelled by $S$, exactly the projectors $q_{\phi,x}$
with $x\in S$ act as one and all others act as zero.  Thus
$H_{0,\Lambda}$ acts as the scalar $|S|$.  This proves
\eqref{eq:spectrum}.  For $|S|=k$ there are $\binom Nk$ choices of
$S$, and every chosen excited factor has dimension $d-1$, proving
\eqref{eq:multiplicity}.

At energy zero, $S=\varnothing$ is the only summand and every local
factor lies in the one-dimensional range of $P_{\phi,x}$.  Its
normalized vector is exactly \eqref{eq:productground}; evenness makes
the product independent of the temporary cluster ordering.  The next
spectral value is one, proving \eqref{eq:gapone}.  No edge or torus
identification enters \eqref{eq:H0}, so the conclusion is independent
of boundary condition.
\end{proof}

\section{Gauge, parity, translations, and degeneracy}

\begin{lemma}[The chiral determinant character is trivial]
\label{lem:det}
For every $g\in\operatorname{Spin}(10)$,
\[
 \det\rho_{16}(g)=1.
\tag{3.1}\label{eq:detone}
\]
\end{lemma}

\begin{proof}
The map $\chi=\det\circ\rho_{16}:
\operatorname{Spin}(10)\to U(1)$ is a continuous group character.
Its differential at the identity is
\[
 d\chi(X)=\Tr(d\rho_{16}(X)).
\]
The trace vanishes on commutators.  Since
$\mathfrak{spin}(10)$ is semisimple,
$[\mathfrak{spin}(10),\mathfrak{spin}(10)]
=\mathfrak{spin}(10)$, so this linear functional vanishes identically.
Thus $\chi$ is locally constant.  The group $\operatorname{Spin}(10)$
is connected, and $\chi(e)=1$, hence $\chi(g)=1$ for every $g$.
\end{proof}

\begin{theorem}[T2: exact symmetries and no hidden degeneracy]
\label{thm:T2}
The Hamiltonian $H_{0,\Lambda}$ has the following properties.
\begin{enumerate}[label=\textup{(\alph*)}]
\item It is invariant under independent local Spin(10) transformations.
Equivalently, for every matter Gauss generator
$G_x^A=d\Gamma_x(T^A)$,
\[
 [H_{0,\Lambda},G_x^A]=0
 \quad(x\in\Lambda,\ A=1,\ldots,45).
\tag{3.2}\label{eq:gauss}
\]
\item It commutes with total fermion parity.  If $\phi$ is constant,
it is invariant under every lattice translation preserving $\Lambda$
and its boundary condition.
\item The finite-volume ground state is a symmetry singlet.  The
infinite-volume product state is the unique frustration-free ground
state, so none of these symmetries is spontaneously broken.
\item On a torus the ground-state degeneracy is exactly one.  In
particular, the model has no topological ground-state degeneracy.
\end{enumerate}
\end{theorem}

\begin{proof}
Second quantization fixes the vacuum.  On the filled vector it acts by
the top exterior power:
\[
 \Gamma(\rho_{16}(g))|F\rangle
 =(\det\rho_{16}(g))|F\rangle=|F\rangle
\]
by Lemma~\ref{lem:det}.  Hence $|\Omega_\phi\rangle$, $P_\phi$, and
$q_\phi$ are invariant.  This is true independently at every cluster,
so differentiation gives \eqref{eq:gauss}.  The forty-five zero
commutators in the frozen finite computation are the matrix twin of
this argument.

Both components of $|\Omega_\phi\rangle$ have even parity because
their particle numbers are $0$ and $16$.  Thus every $P_{\phi,x}$ and
$q_{\phi,x}$ is even and $H_0$ commutes with total parity.  A
translation merely permutes identical terms when $\phi$ is constant,
which proves translation invariance.

The product ground vector is fixed by all the stated symmetries.  For
the infinite system, let $\omega$ be any frustration-free state.
Then $\omega(q_{\phi,x})=0$ for every $x$.  Positivity and
$q_{\phi,x}\geq0$ imply that the one-site reduced state is supported
on $\ker q_{\phi,x}=\C|\Omega_\phi\rangle$.  A state with a pure
one-site marginal factorizes across that site; iterating over every
finite region gives
\[
 \omega|_A
 =\bigotimes_{x\in A}
 |\Omega_\phi\rangle\langle\Omega_\phi|
\]
for every finite $A$.  Local restrictions determine a quasi-local
state, so $\omega$ is the unique product state.  There are therefore
no distinct infinite-volume ground states related by a symmetry and
no spontaneous symmetry breaking.

Finally, Theorem~\ref{thm:T1} applies unchanged on a torus and gives a
one-dimensional ground space.  More concretely, all local constraints
$q_{\phi,x}=0$ force the same rank-one factor independently at every
site; no noncontractible cycle supplies an additional label.  This
proves the torus claim directly from the commuting-projector
structure.
\end{proof}

\section{The classical stability theorem, stated in full}

\subsection{Local-TQO and local gap}

For a finite region $B\subseteq\Lambda$, define
\[
 H_B=\sum_{x\in B}q_{\phi,x},\qquad
 P_B=\prod_{x\in B}P_{\phi,x}.
\tag{4.1}\label{eq:localHB}
\]
As an operator on the full system, $P_B$ is the identity outside
$B$.  If $A\subset B$ and $O_A$ is supported on $A$, then
\[
 P_BO_AP_B
 =\ip{\Omega_{A,\phi}}{O_A\Omega_{A,\phi}}P_B.
\tag{4.2}\label{eq:exactLTQO}
\]
Thus local ground states are exactly indistinguishable in the
interior: the Local-TQO error function is identically zero.  Theorem
\ref{thm:T1} applied to $B$ also gives
\[
 \gap(H_B)=1
\tag{4.3}\label{eq:localgap}
\]
for every nonempty $B$.

\begin{definition}[Finite-range perturbation strength]
\label{def:pert}
An even finite-range perturbation is a sum
\[
 V_\Lambda=\sum_{X\subseteq\Lambda}V_X,\qquad
 V_X=V_X^*,\qquad |V_X|=0,
\tag{4.4}\label{eq:pert}
\]
for which $V_X=0$ when $\operatorname{diam}X>R$ and
\[
 \sup_{x\in\Lambda}\sum_{X\ni x}\norm{V_X}\leq J_V
\tag{4.5}\label{eq:interactionnorm}
\]
uniformly in $\Lambda$.  Gauge covariance means that the complete
local term, including any prescribed parallel transporter, is
Spin(10)-invariant.  Evenness makes disjoint interaction terms commute
in the graded sense required by the local stability argument.
\end{definition}

\begin{theorem}[Michalakis--Zwolak stability theorem, full form used
here]
\label{thm:MZ}
Let $H_0$ be a finite-range frustration-free Hamiltonian on finite
subsets of $\mathbb Z^D$, with a volume-independent spectral gap
$\gamma>0$.  For every ball $b_u(r)$ let $P_{b_u(r)}$ be the ground
projection of the restricted Hamiltonian.  Assume:
\begin{enumerate}[label=\textup{(\roman*)}]
\item \emph{Local-TQO.}  For every ball $A=b_u(r)$, every bounded
operator $O_A$ supported in $A$, and
$A(\ell)=b_u(r+\ell)$,
\[
 \left\|
 P_{A(\ell)}O_AP_{A(\ell)}
 -\frac{\Tr(P_{A(\ell)}O_A)}
 {\Tr P_{A(\ell)}}P_{A(\ell)}
 \right\|
 \leq\norm{O_A}\Delta_0(\ell),
\tag{4.6}\label{eq:MZLTQO}
\]
where $\Delta_0(\ell)$ decays sufficiently rapidly and the admissible
topological-order length grows as a positive power of the system
diameter.
\item \emph{Local gap.}  Every restricted Hamiltonian on $b_u(r)$ has
no spectrum in $(0,\gamma(r))$, where $\gamma(r)>0$ decays at most
polynomially in $r$.
\item \emph{Quasi-local perturbation.}  The perturbation is a
$(J,f_0)$ interaction: it is a sum of bounded terms centred at lattice
sites whose range-$r$ norms are bounded by $Jf_0(r)$, with $f_0$
decaying rapidly enough for a Lieb--Robinson bound.
\end{enumerate}
Then there exist constants $J_0>0$ and $L_0\geq2$, independent of the
finite volume, such that whenever $J\leq J_0$ and the system diameter
$L\geq L_0$,
\[
 \gap(H_0+V)\geq\frac{\gamma}{2}.
\tag{4.7}\label{eq:MZgap}
\]
In the proof, after quasi-local spectral-flow block diagonalization,
the transformed interaction $W$ obeys
\[
 |\ip{\psi}{W\psi}|
 \leq cJ\,\ip{\psi}{H_0\psi},
\tag{4.8}\label{eq:MZrelative}
\]
where one may take
\[
 c=C_D\sum_{k=1}^{M}
 \frac{r_k^D}{\gamma(r_k)}
 \sum_{r_{k-1}<r\leq r_k}w(r),
\qquad
 0=r_0<r_1\leq\cdots\leq r_M=L,
\tag{4.9}\label{eq:MZc}
\]
$C_D$ is the volume of the unit ball in $\mathbb Z^D$, and $w$ is the
rapid-decay function of the transformed interaction.  The final
bootstrap may choose
\[
 J_0=\frac1{3c}.
\tag{4.10}\label{eq:MZJ0}
\]
The constants $w,c,J_0$ depend on the interaction-decay and
spectral-flow estimates, but not on volume.
\end{theorem}

This is Theorem~1 and Proposition~2 of
S.~Michalakis and J.~P.~Zwolak,
\emph{Stability of Frustration-Free Hamiltonians},
Communications in Mathematical Physics \textbf{322} (2013), 277--302.
It extends the commuting-projector theorem of Bravyi, Hastings, and
Michalakis.  We have included the hypotheses, conclusion, and the
relative-bound constant used in their proof so that no numerical
threshold is hidden.

\section{Volume-independent perturbative stability}

\begin{theorem}[T3: proved volume-independent stability]
\label{thm:T3}
Let $\Lambda$ range over finite subsets or tori of a fixed
$D$-dimensional bounded-degree lattice.  Let $H_{0,\Lambda}$ be the
DET16 Hamiltonian \eqref{eq:H0}, and let $V_\Lambda$ be an even
finite-range interaction satisfying Definition~\ref{def:pert}
uniformly in volume.  Then there exists $t_*>0$, depending only on
$D$, the interaction range, and the local interaction-norm bound, but
not on $\Lambda$, such that
\[
 \gap(H_{0,\Lambda}+tV_\Lambda)\geq\frac12
\tag{5.1}\label{eq:stablehalf}
\]
for every sufficiently large volume and $|t|\leq t_*$.  After
decreasing $t_*$ once to cover the finitely many smaller volumes, the
same statement holds for every finite $\Lambda$.
\end{theorem}

\begin{proof}
Every local term $q_{\phi,x}$ is positive and annihilates the product
state, so $H_0$ is frustration free.  Theorem~\ref{thm:T1} gives
global gap $\gamma=1$.  Equation \eqref{eq:exactLTQO} verifies the full
Local-TQO condition \eqref{eq:MZLTQO} with
\[
 \Delta_0(\ell)=0
\]
at every scale, not merely asymptotically.  Equation
\eqref{eq:localgap} verifies the Local-Gap condition with
$\gamma(r)=1$.  A finite-range interaction has a compactly supported,
hence rapidly decaying, function $f_0$.  The local on-cluster Hilbert
space is finite dimensional.  Finally, all perturbation terms are
even, so spacelike-separated local terms commute rather than
anticommute; the locality and Lieb--Robinson estimates used in
Theorem~\ref{thm:MZ} apply to this even CAR interaction.

Apply Theorem~\ref{thm:MZ} to the perturbation $tV$.  Its strength is
$J=|t|J_V$, so one may set $t_*=J_0/J_V$.  This gives
\eqref{eq:stablehalf} for $L\geq L_0$.  There are only finitely many
lattice shapes below $L_0$ up to the fixed local geometry.  Each has a
simple isolated ground eigenvalue and gap one at $t=0$; finite-matrix
eigenvalues are continuous in $t$.  Taking the minimum of their
positive continuity radii and the large-volume value gives one
positive threshold for all finite volumes.
\end{proof}

\section{Dynamical finite gauge links: Lie--Schwinger gap theorem}
\label{sec:dynlinks}

Theorem~\ref{thm:T3} applies to prescribed parallel transporters but does
not itself add fluctuating link Hilbert spaces.  We therefore audit the
finite-group dynamical-link extension first against the local
Schrieffer--Wolff construction of Bravyi, DiVincenzo, and
Loss~\cite{BDL2011}, and then against the gap theorem proved by
Fröhlich and Pizzo~\cite{FP2020}.

Let each oriented open-chain link carry the electric basis of
$\C[\mathbb Z_N]$, $N=2$ or $4$, and let $U_\ell$ shift its electric
flux.  On the unconstrained ambient tensor product put
\[
 H_{0,\mathrm{dyn}}
 =\sum_x(\id-P_{\phi,x})+h_E\sum_\ell E_\ell^2,\qquad
 V=\sum_{\ell=(x,x+1)}K_\ell ,
\tag{5.2}\label{eq:dynH0}
\]
where
\[
 K_\ell=\sum_{\alpha=1}^{m}
 \left(c^\dagger_{x,\alpha}U_\ell c_{x+1,\alpha}
 +c^\dagger_{x+1,\alpha}U_\ell^\dagger c_{x,\alpha}\right).
\tag{5.3}\label{eq:dynhop}
\]
Here $m=16$ for DET16 and $m=4$ in the numeric analogue.  Assigning each
outgoing link to its left matter cell makes \eqref{eq:dynH0} a sum of
one-cell terms and \eqref{eq:dynhop} a two-cell interaction.  Its product
ground projector is the tensor product of the DET projectors and the
zero-electric-flux link projectors, and its local gap is
\[
 \Delta=\min(1,h_E).
\tag{5.4}\label{eq:dyndelta}
\]
Moreover $\norm{K_\ell}=m$, so the BDL interaction strength obeys
\[
 J=\max_x\sum_{\ell\ni x}\norm{K_\ell}\leq2m
\tag{5.5}\label{eq:dynJ}
\]
uniformly in chain length.

\begin{proposition}[BDL hypothesis repair and remaining obstruction]
\label{prop:BDLaudit}
For finite $\mathbb Z_2$ or $\mathbb Z_4$ links, the Gauss-restricted
Hilbert space is not a tensor product and hence is not directly in the
setup of Section~4 of Bravyi--DiVincenzo--Loss.  The ambient-cell
formulation \eqref{eq:dynH0}--\eqref{eq:dynJ} repairs this locality
hypothesis, and gauge invariance makes every Gauss sector reducing.
Nevertheless the cited BDL result does not supply the requested
volume-independent spectral-gap theorem or a numerical threshold.
\end{proposition}

\begin{proof}
The local Gauss constraints correlate a matter cell with its incident
links, so the physical subspace does not factor into local Hilbert spaces.
On the ambient space, by contrast, the left-link cell assignment gives
exactly the tensor product, one-cell unperturbed Hamiltonian, product
ground space, and two-cell perturbation assumed in BDL Section~4.1.
The determinant changes charge by sixteen and the dressed hopping
transports one unit of link flux with the matter charge, so all terms
commute with every Gauss generator.  Restriction to the physical sector
therefore preserves the spectrum originating from its invariant
subspace.  Adding a Gauss penalty is unnecessary; it would produce
overlapping commuting local terms and would no longer have BDL's
one-site unperturbed form.

For a second-order truncation $n=2$, BDL Theorem~3 normalizes
$\norm{V}_1=1$ and obtains
\[
 \epsilon_c
 =\min\left\{\frac{\Delta\beta}{4n^2},
              \frac{\Delta}{32\alpha}\right\},
\tag{5.6}\label{eq:BDLeps}
\]
where $\alpha$ and $\beta$ arise from big-$O$ estimates in their
Lemma~4.2 and are not evaluated numerically.  Writing
$\epsilon=tJ$ and using \eqref{eq:dynJ} gives only the symbolic
conditions
\[
\begin{aligned}
 m=4:\quad
 |t|&<\Delta\min\left\{\frac{\beta}{128},
                  \frac1{256\alpha}\right\},\\
 m=16:\quad
 |t|&<\Delta\min\left\{\frac{\beta}{512},
                  \frac1{1024\alpha}\right\}.
\end{aligned}
\tag{5.7}\label{eq:BDLsymbolic}
\]
No concrete number follows without values of $\alpha,\beta$.
More decisively, BDL Theorem~3 controls the \emph{ground-state energy}
with an extensive truncation error
$O(|\Lambda||\epsilon|^{n+1})$; it does not control the low-sector
spectral gap.  The paragraph following their Lemma~4.1 explicitly states
that the lemma says nothing about that gap.  Consequently
\eqref{eq:BDLsymbolic} cannot be promoted to a uniform lower bound on
$\gap(H_{0,\mathrm{dyn}}+tV)$.
\end{proof}

\begin{theorem}[Fröhlich--Pizzo Lie--Schwinger theorem, full form used
here]
\label{thm:FP}
Let
$\Hh^{(L)}=\bigotimes_{j=1}^{L}\Hh_j$ with
$\Hh_j\simeq\C^M$, where $M<\infty$ is independent of $L$.
Suppose the identical one-site operators $H_j\geq0$ have a
one-dimensional kernel $\C\Omega_j$ and satisfy
\[
 H_j|_{\{\C\Omega_j\}^{\perp}}\geq\id.
\tag{5.8}\label{eq:FPlocal}
\]
Let
\[
 K_L(\tau)=\sum_{j=1}^{L}H_j+
 \tau\sum_{\substack{I_{k;i}\subset\{1,\ldots,L\}\\
                     k\leq\bar k}}V_{I_{k;i}},
\tag{5.9}\label{eq:FPK}
\]
where $\bar k<\infty$ is fixed independently of $L$, every
$V_{I_{k;i}}$ is self-adjoint and supported on its indicated interval,
and $\norm{V_{I_{k;i}}}\leq1$.  Then there is
$\tau_{\rm FP}>0$, independent of $L$, such that for every real
$|\tau|<\tau_{\rm FP}$ and every finite $L$, $K_L(\tau)$ has a unique
ground state and
\[
 \gap K_L(\tau)\geq\frac12.
\tag{5.10}\label{eq:FPgap}
\]
\end{theorem}

This is Theorem~3.5 of
J.~Fröhlich and A.~Pizzo,
\emph{Lie--Schwinger Block-Diagonalization and Gapped Quantum Chains},
Communications in Mathematical Physics \textbf{375} (2020),
2039--2069~\cite{FP2020}.  Its conclusion is a spectral-gap estimate,
not a ground-energy truncation estimate, and is uniform in chain
length.

\begin{proposition}[Lie--Schwinger hypothesis audit for finite gauge
links]
\label{prop:FPaudit}
For $N_g\in\{2,4\}$, $h_E>0$, $m\in\{4,16\}$, and constant $\phi$, the
open-chain Hamiltonian \eqref{eq:dynH0}--\eqref{eq:dynhop} satisfies
every hypothesis of Theorem~\ref{thm:FP} on an ambient tensor product
after adjoining one decoupled endpoint link.  With
\[
 \Delta=\min(1,h_E),\qquad
 \tau=\frac{tm}{\Delta},
\tag{5.11}\label{eq:FPscale}
\]
therefore
\[
 |t|<t_{\rm FP}^{(m)}(h_E)
 :=\frac{\Delta}{m}\tau_{\rm FP}
 \quad\Longrightarrow\quad
 \gap\!\left(
 H_{0,\mathrm{dyn}}+tV
 \right)\big|_{\Hh_{\rm phys}}
 \geq\frac{\Delta}{2},
\tag{5.12}\label{eq:FPwindow}
\]
and the physical Gauss sector has a unique ground state.  The bound is
uniform in the number of clusters.
\end{proposition}

\begin{proof}
The audit is hypothesis by hypothesis:
\begin{table}[htbp]
\centering
\small
\begin{tabularx}{\textwidth}{@{}p{0.25\textwidth}X p{0.10\textwidth}@{}}
\toprule
\textbf{FP hypothesis}&\textbf{DET chain realization}&\textbf{Result}\\
\midrule
Finite tensor-product cells&
$\Hh_x=\F(\C^m)\otimes\C[\mathbb Z_{N_g}]$, of dimension
$2^mN_g$ independent of chain length.&Pass\\
Identical one-site $H_x$&
Put the outgoing link in cell $x$ and append one decoupled final link,
so every cell carries
$q_{\phi,x}+h_EE_x^2$.&Pass\\
Rank-one local vacuum&
$P_{\phi,x}$ and the zero-electric-flux projector both have rank one;
their tensor product is the unique cell vacuum.&Pass\\
Normalized local gap&
The cell gap is $\Delta=\min(1,h_E)>0$; dividing $H_x$ by $\Delta$
gives \eqref{eq:FPlocal}.&Pass\\
Fixed finite interval range&
$K_{x,x+1}$ acts on two adjacent cells, hence $\bar k=1$ independently
of volume.&Pass\\
Ordinary tensor locality&
In an ordered Jordan--Wigner representation the adjacent strings
cancel; the even bilinear hopping, including its left-cell link
operator, acts on exactly two cells.&Pass\\
Bounded self-adjoint interaction&
$K_{x,x+1}=K_{x,x+1}^{\dagger}$ and
$\norm{K_{x,x+1}}=m$; thus
$V_{x,x+1}=K_{x,x+1}/m$ has norm one.&Pass\\
Real small coupling&
After factoring out $\Delta$, the FP coupling is precisely
$\tau=tm/\Delta$.&Pass\\
Gauss restriction&
This is not an FP premise: all terms commute with every Gauss
generator, and the ambient vacuum is physical because
$m=0\pmod {N_g}$ with zero link flux.&Pass\\
\bottomrule
\end{tabularx}
\caption{Fröhlich--Pizzo hypotheses audited against the finite-link
DET chain.}
\label{tab:FPaudit}
\end{table}
The appended endpoint link is decoupled and has a unique zero-flux
ground state.  The enlarged Hamiltonian is therefore the original
ambient Hamiltonian tensored with this one-link factor; a lower bound
on the enlarged gap is also a lower bound on the original ambient gap.

The normalized Hamiltonian is exactly \eqref{eq:FPK} with
$\bar k=1$, so Theorem~\ref{thm:FP} gives a unique ambient ground state
and normalized gap at least $1/2$.  Every Gauss generator commutes with
the Hamiltonian.  At $t=0$ the unique ambient vacuum is in the trivial
Gauss sector, since both $m=4$ and $m=16$ vanish modulo $2$ and $4$.
Uniqueness and the discrete gauge eigenvalues keep the continued ground
state in that sector.  The physical spectrum is a subset of the ambient
spectrum and contains its ground state, so restriction cannot reduce
the excitation gap.  Restoring the factor $\Delta$ proves
\eqref{eq:FPwindow}.
\end{proof}

\subsection{What is and is not explicit in the FP window}

The Lie--Schwinger proof contains one explicit series constant:
\[
 \frac{e^{8a}-8a-1}{a}+e^{8a}-1=1,\qquad
 a=0.023320199830777650\ldots .
\tag{5.14}\label{eq:FPa}
\]
Appendix~A of \cite{FP2020} uses the explicit convergence restriction
\[
 |\tau|<\frac a{16}
 =0.001457512489423603\ldots .
\]
Its local-gap estimate also has the explicit sufficient inequality
\[
 8|\tau|+16|\tau|
 \sum_{\ell=3}^{\infty}
 \frac{|\tau|^{(\ell-2)/3}}{\ell}\leq\frac12,
\tag{5.15}\label{eq:FPgapcondition}
\]
whose positive boundary is
$\tau_{\rm gap}=0.046907169193748369\ldots$.
However, the induction closing Theorem~3.4 of \cite{FP2020} introduces
unevaluated universal constants $C,C'$ and repeatedly requires
``sufficiently small'' $|\tau|$.  A scrupulous parameterization is
therefore
\[
 \tau_{\rm FP}
 =\min\left\{\frac a{16},\tau_{\rm gap},\tau_{\rm ind}\right\}>0,
\tag{5.16}\label{eq:FPparametric}
\]
where $\tau_{\rm ind}>0$ denotes those volume-independent but
numerically unevaluated induction restrictions.  Thus the two concrete
model families have the cited, parametric windows
\[
\begin{array}{c|cc}
\toprule
&m=4\ \text{analogue}&m=16\ \text{DET16}\\
\midrule
h_E=\tfrac12&|t|<\tau_{\rm FP}/8&|t|<\tau_{\rm FP}/32\\
h_E\geq1&|t|<\tau_{\rm FP}/4&|t|<\tau_{\rm FP}/16\\
\bottomrule
\end{array}
\tag{5.17}\label{eq:FPwindows}
\]
No numerical value of $t_{\rm FP}^{(m)}$ follows unless
$\tau_{\rm ind}$ is evaluated.

The deterministic companion
\[
 \texttt{experiments/tfpt-discovery/
 det16\_dynamical\_links\_probe.py}
\]
passes all $149$ checks.  It verifies exact Gauss invariance, commuting
unperturbed terms, rank-one local vacua, the norms
\eqref{eq:dynJ}, the FP range and normalization hypotheses, the two
explicit constants \eqref{eq:FPa}--\eqref{eq:FPgapcondition}, the
second-order effective Hamiltonian, and fourth-order finite-volume
residual scaling.  A long-range mutant supported on the two endpoints
has diameter $L-1$ and visibly fails the fixed-$\bar k$ premise.
The smallest scanned gap remains $0.728854416$ through $|t|=0.4$ on
two and three four-mode clusters.  At $h_E\geq1$, that scan point has
normalized coupling $\tau=1.6$ for $m=4$ and $\tau=6.4$ for $m=16$,
both far outside the cited window; the scan is evidence against a
finite-volume closing there, not an extension of
\eqref{eq:FPwindow}.  The finite $\mathbb Z_2/\mathbb Z_4$ chain item is
therefore cited-verified with the parametric FP window, not with an
invented numerical threshold.

\section{The proposed explicit count and why it is conditional}

\begin{proposition}[Elementary relative-form gap bound]
\label{prop:relative}
Let $P=|\Omega_{\Lambda,\phi}\rangle
\langle\Omega_{\Lambda,\phi}|$ and $Q=\id-P$.  Suppose a self-adjoint
perturbation $V$ satisfies
\[
 PVQ=QVP=0,\qquad PVP=v_0P,
\tag{6.1}\label{eq:block}
\]
and
\[
 \left|
 \ip{\psi}{(V-v_0\id)\psi}
 \right|
 \leq c\,\ip{\psi}{H_0\psi}
 \quad\text{for every }\psi\in Q\Hh.
\tag{6.2}\label{eq:formbound}
\]
Then the ground vector remains $|\Omega_{\Lambda,\phi}\rangle$, and
\[
 \gap(H_0+tV)\geq1-c|t|
 \geq1-2c|t|.
\tag{6.3}\label{eq:linearbound}
\]
In particular, the weaker displayed bound is positive for
$|t|<t_{0,\mathrm{rel}}:=1/(2c)$.
\end{proposition}

\begin{proof}
Condition \eqref{eq:block} makes $P\Hh$ and $Q\Hh$ reducing subspaces.
The energy on $P\Hh$ is $tv_0$.  For normalized $\psi\in Q\Hh$,
Theorem~\ref{thm:T1} gives
$\ip{\psi}{H_0\psi}\geq1$, and
\[
\begin{aligned}
 \ip{\psi}{(H_0+tV-tv_0\id)\psi}
 &\geq
 (1-c|t|)\ip{\psi}{H_0\psi}\\
 &\geq1-c|t|.
\end{aligned}
\]
The min--max principle proves the first inequality in
\eqref{eq:linearbound}; the second is a weakening for $c|t|\geq0$.
\end{proof}

\subsection{The norm count}

For unit-amplitude nearest-neighbour hopping across one bond,
\[
 K_{xy}=\sum_{\alpha,\beta=1}^{16}
 \left(a_{x,\alpha}^\dagger U_{\alpha\beta}a_{y,\beta}
 +a_{y,\beta}^\dagger U_{\alpha\beta}^*a_{x,\alpha}\right),
\tag{6.4}\label{eq:hopping}
\]
with $U$ unitary, the one-particle bond matrix has sixteen eigenvalues
$+1$ and sixteen eigenvalues $-1$.  Its fermionic second quantization
therefore has
\[
 \norm{K_{xy}}=16.
\tag{6.5}\label{eq:bondnorm}
\]
If the lattice coordination is $z$, at most $z$ such bonds touch one
cluster.  The advertised count is consequently
\[
 c_{\mathrm{count}}=16z,\qquad
 t_{0,\mathrm{count}}=\frac1{32z}.
\tag{6.6}\label{eq:count}
\]
On a three-dimensional cubic spatial slice, $z=6$, so
\[
 c_{\mathrm{count}}=96,\qquad
 t_{0,\mathrm{count}}=\frac1{192}
 =0.005208333\ldots .
\tag{6.7}\label{eq:count3d}
\]

This count controls the norm incident on one site.  It does not prove
\eqref{eq:block} or \eqref{eq:formbound}.  In fact ordinary hopping
creates a pair of defective clusters from components of the product
ground state, so
\[
 QK P\ne0.
\tag{6.8}\label{eq:notblock}
\]
The numeric twin measures
$\norm{QKP}=\sqrt2$ on the two-cluster four-mode analogue.  Thus
Proposition~\ref{prop:relative} cannot be applied to
\eqref{eq:hopping}, and \eqref{eq:count3d} is not the explicit
$t_*$ of Theorem~\ref{thm:T3}.  A global Weyl estimate would only give
$1-2|t|\norm V$, whose right side scales with volume and is unsuitable.

\begin{center}
\setlength{\fboxsep}{9pt}
\fcolorbox{orange!75!black}{orange!4}{
\begin{minipage}{0.91\textwidth}
\textbf{Uncertified explicit-threshold box.}
The number $t_{0,\mathrm{count}}=1/192$ is the value obtained from the
requested local norm count on a cubic spatial slice.  It is a valid
threshold for the \emph{conditional} relative-form proposition if
\eqref{eq:block}--\eqref{eq:formbound} are separately proved with
$c=96$.  Ordinary hopping fails \eqref{eq:block}, visibly and
numerically.  The applicable Michalakis--Zwolak theorem instead gives
$t_*=J_0/J_V>0$ with $J_0=1/(3c_{\rm MZ})$, where $c_{\rm MZ}$ contains
the quasi-local spectral-flow decay $w$ in \eqref{eq:MZc}.  No value of
$c_{\rm MZ}$ has been derived for this model, so replacing it by $96$
would be an unjustified theorem claim.
\end{minipage}}
\end{center}

\section{Numeric twin}

The companion probe is
\[
 \texttt{experiments/tfpt-discovery/det16\_gap\_stability\_probe.py}.
\]
It uses the exact four-mode determinant projectors and open-chain
hopping machinery of the frozen \texttt{v1002} route.  All $19$ checks
pass.  The direct comparison is:
\[
\begin{array}{c|rrrr}
\toprule
\text{clusters}&t=0.05&t=0.10&t=0.20&t=0.30\\
\midrule
2&0.973726995&0.946980790&0.902548402&0.837138848\\
3&0.962836264&0.924929608&0.860366685&0.764340648\\
4&0.957577627&0.914546582&0.842484121&0.726091276\\
\bottomrule
\end{array}
\tag{7.1}\label{eq:numeric}
\]
Every entry lies above $1-2c_{\rm count}|t|$ with
$c_{\rm count}=4$ for the two-cluster endpoint pair and $8$ for the
three--four cluster chains.  Except at the smallest $t$, this linear
line is already negative and hence not tight.

The scan uses step $0.05$ for two and three clusters and step $0.10$
for four clusters on $0\leq t\leq4$.  It finds no finite-volume gap
closing.  The minimum scanned gaps are
\[
\begin{array}{c|cc}
\toprule
\text{clusters}&\min\gap&\text{minimizing }t\\
\midrule
2&0.764490286&0.55\\
3&0.472500856&4.00\\
4&0.549772255&0.60\\
\bottomrule
\end{array}
\tag{7.2}\label{eq:scan}
\]
Consequently the requested finite-chain ``actual closing''
$t_c$ is not located because it does not occur in the declared scan:
\[
 t_c^{\rm scan}>4,\qquad
 \frac{t_{0,\mathrm{count}}^{\rm toy}}{t_c^{\rm scan}}
 <\frac{1/16}{4}=0.015625.
\tag{7.3}\label{eq:margin}
\]
This is an upper bound on a ratio, not an invented critical point.

The positive and negative controls expose the theorem hypotheses:
\[
\begin{array}{ll}
\text{Spin(4)-covariant flavor-diagonal hopping:}
 &\max_A\norm{[K,G^A]}=0,\\
\text{flavor-selective mutant:}
 &\max_A\norm{[K_{\rm mut},G^A]}=1,\\
\text{$1/3$ normalization mutant:}
 &\norm{\widetilde P^2-\widetilde P}=1/9,\quad
 \Tr\widetilde P=2/3.
\end{array}
\tag{7.4}\label{eq:mutants}
\]
Thus one mutant visibly leaves the covariant perturbation class and the
other destroys the projector/frustration-free premise.

\section{Combined strongest theorem}

\begin{theorem}[Volume-independent DET16 mirror-gap theorem]
\label{thm:main}
For the matter-only DET16 cluster model
\eqref{eq:cluster}--\eqref{eq:H0} on any finite bounded-degree lattice:
\begin{enumerate}[label=\textup{(\alph*)}]
\item the projectors are mutually commuting even CAR operators;
$H_0$ has exact spectrum $\{0,\ldots,|\Lambda|\}$, multiplicities
\eqref{eq:multiplicity}, unique product ground state
\eqref{eq:productground}, and gap one in every volume;
\item $H_0$ is locally Spin(10)-Gauss invariant, parity invariant, and
translation invariant for constant $\phi$; its unique product ground
state has neither spontaneous symmetry breaking nor torus/topological
degeneracy;
\item every uniformly bounded finite-range even perturbation has a
volume-independent stability interval $|t|\leq t_*$ on which the gap
is at least $1/2$, with $t_*>0$ supplied by the
Michalakis--Zwolak theorem after the exact checks
$\Delta_0=0$ and $\gamma(r)=1$;
\item for finite $\mathbb Z_2/\mathbb Z_4$ links on an open chain,
$h_E>0$, and gauge-covariant nearest-neighbour hopping, the physical
Gauss sector has a unique ground state and volume-independent gap at
least $\Delta/2$ throughout the cited window
$|t|<\Delta\tau_{\rm FP}/16$, with
$\Delta=\min(1,h_E)$ and $\tau_{\rm FP}>0$ parameterized by
\eqref{eq:FPparametric};
\item if, additionally, a perturbation obeys the block-diagonal
relative-form hypotheses \eqref{eq:block}--\eqref{eq:formbound}, then
the explicit elementary estimate \eqref{eq:linearbound} holds.
\end{enumerate}
No numerical identification of $t_*$ with
$t_{0,\mathrm{count}}=1/(32z)$ is asserted.
\end{theorem}

\begin{proof}
Part (a) is Theorem~\ref{thm:T1} and Lemma~\ref{lem:order}.
Part (b) is Theorem~\ref{thm:T2}.  Part (c) is
Theorem~\ref{thm:T3}, whose proof verifies every hypothesis of the
quoted stability theorem against the concrete product model.  Part
(d) is Proposition~\ref{prop:FPaudit}, including the ambient-to-physical
gap restriction.  Part (e) is Proposition~\ref{prop:relative}.  The
final sentence follows from \eqref{eq:notblock} and the
explicit-threshold box.
\end{proof}

\section{Physical spectator and remaining boundary}
\label{sec:boundary}

\subsection{Exact spectator statement}

Let $\Hh_{\rm wall}$ carry a wall Hamiltonian $H_{\rm wall}$ and let
$\Hh_{\rm mir}$ carry $H_{0,\Lambda}$.  If the DET16 terms act as the
identity on wall modes, then
\[
 H_{\rm total}(0)
 =H_{\rm wall}\otimes\id_{\rm mir}
 +\id_{\rm wall}\otimes H_{0,\Lambda}.
\tag{9.1}\label{eq:spectator}
\]
Hence
\[
 \operatorname{spec}H_{\rm total}(0)
 =\{E_{\rm wall}+k:
 E_{\rm wall}\in\operatorname{spec}H_{\rm wall},\
 k=0,\ldots,|\Lambda|\}.
\tag{9.2}\label{eq:spectralsum}
\]
The $k=0$ copy is exactly the original wall spectrum, with every
multiplicity preserved.

For $t\ne0$, ``does not mix'' must mean that no interaction term
contains both wall and mirror operators.  Exact preservation, rather
than merely perturbative continuity, further requires
\[
 V=\id_{\rm wall}\otimes V_{\rm mir}.
\tag{9.3}\label{eq:nomix}
\]
Then the total Hamiltonian remains a tensor sum, and the stable mirror
ground band adds one common scalar energy to the wall spectrum.  If a
wall-only term is allowed it changes the wall spectrum on its own; if a
wall--mirror coupling is allowed, \eqref{eq:spectralsum} is no longer
exact and a separate relative/local stability estimate is required.

\begin{center}
\setlength{\fboxsep}{10pt}
\fcolorbox{red!65!black}{red!3}{
\begin{minipage}{0.91\textwidth}
\textbf{Remaining boundary for
\texttt{CHIRAL4D.MIRROR.DET16.01}.}

\smallskip
\textbf{Compact/continuous dynamical gauge fields.}
Section~\ref{sec:dynlinks} closes the finite
$\mathbb Z_2/\mathbb Z_4$ open-chain item by
Theorem~\ref{thm:FP} and Proposition~\ref{prop:FPaudit}; it is no longer
part of this boundary box.  Compact-group rotor links have
infinite-dimensional local Hilbert spaces and therefore violate the
finite-$M$ hypothesis of the bounded-interaction FP theorem quoted here.
They require the form-bounded extension with its domain hypotheses, or
truncation-uniform estimates, neither of which is audited in this
memorandum.

\smallskip
\textbf{Domain wall.}
At $t=0$ the spectator theorem
\eqref{eq:spectator}--\eqref{eq:spectralsum} is exact.  For
$|t|<t_*$, exact preservation of the wall spectrum requires the
factorization \eqref{eq:nomix}.  Gauge covariance or fermion parity
alone does not forbid wall--mirror bilinears, so the microscopic
domain-wall construction must prove their absence or bound them.

\smallskip
\textbf{Three spatial dimensions.}
The count $z=6$ and $1/192$ assumes a cubic three-dimensional spatial
slice, but no lattice domain-wall Hamiltonian, chirality localization,
continuum limit, or Lorentz-covariant $3+1$-dimensional theory is
constructed here.  The theorem is a uniform quantum-lattice statement
for the mirror matter sector, not yet the physical four-dimensional
placement.

\smallskip
\textbf{Explicit hopping threshold.}
The applicable stability theorem proves $t_*>0$ and gap at least
$1/2$, but its constant contains the spectral-flow decay
$c_{\rm MZ}$.  The simple value $1/192$ is conditional and cannot be
promoted to the hopping theorem because $QKP\ne0$.  Deriving a
numerical $c_{\rm MZ}$, or proving a different explicit local
block-diagonalization/Knabe estimate, remains open.
\end{minipage}}
\end{center}

\section*{Conclusion}
\addcontentsline{toc}{section}{Conclusion}

T1 and T2 are elementary, exact, and volume independent.  T3 is a
genuine thermodynamic stability statement: the product projector has
zero Local-TQO error and unit local gap, so the full
Michalakis--Zwolak theorem applies to bounded finite-range even
perturbations.  Its threshold is uniform but not numerically evaluated.
The requested count $1/(32z)$ is retained only in the precise
conditional proposition it actually satisfies; the numeric twin both
passes the comparison and detects the missing block-diagonal premise.
The dynamical-link audit repairs the tensor-product endpoint exactly and
applies the Fröhlich--Pizzo theorem, which does control the spectral gap
uniformly in chain length.  This moves finite
$\mathbb Z_2/\mathbb Z_4$ open chains to cited-verified status with the
parametric window $|t|<\Delta\tau_{\rm FP}/m$; it does not manufacture a
numerical value for the paper's implicit induction constant.  Accordingly,
the red box now retains only compact/continuous links, domain-wall
mixing, and the $3+1$-dimensional placement, preventing the chain theorem
from being mistaken for the complete dynamical-gauge contract.

\begin{thebibliography}{9}

\bibitem{BHM2010}
S.~Bravyi, M.~B.~Hastings, and S.~Michalakis,
\emph{Topological quantum order: stability under local perturbations},
J.\ Math.\ Phys.\ \textbf{51} (2010), 093512,
\href{https://doi.org/10.1063/1.3490195}
{doi:10.1063/1.3490195}.

\bibitem{MZ2013}
S.~Michalakis and J.~P.~Zwolak,
\emph{Stability of Frustration-Free Hamiltonians},
Commun.\ Math.\ Phys.\ \textbf{322} (2013), 277--302,
\href{https://doi.org/10.1007/s00220-013-1762-6}
{doi:10.1007/s00220-013-1762-6}.

\bibitem{BDL2011}
S.~Bravyi, D.~P.~DiVincenzo, and D.~Loss,
\emph{Schrieffer--Wolff transformation for quantum many-body systems},
Ann.\ Phys.\ \textbf{326} (2011), 2793--2826,
\href{https://doi.org/10.1016/j.aop.2011.06.004}
{doi:10.1016/j.aop.2011.06.004}.

\bibitem{FP2020}
J.~Fröhlich and A.~Pizzo,
\emph{Lie--Schwinger Block-Diagonalization and Gapped Quantum Chains},
Commun.\ Math.\ Phys.\ \textbf{375} (2020), 2039--2069,
\href{https://doi.org/10.1007/s00220-019-03613-2}
{doi:10.1007/s00220-019-03613-2};
\href{https://arxiv.org/abs/1812.02457}
{arXiv:1812.02457}.

\end{thebibliography}

\end{document}
